\documentclass[12pt]{article}
\usepackage[UTF8]{inputenc}
\usepackage{thaisupp}
\usepackage{enumitem}
\begin{document}
\title{ของแข็งและความยืดหยุ่น}
\maketitle
\section*{In-class Questions}
\begin{enumerate}[label=\arabic*.]
\item ความเค้น (Stress) คืออะไร \hfill \textbf{\small (Main)}
\begin{enumerate}[label=)]
\item ก) แรงที่กระทำต่อวัตถุ
\item ข) แรงต่อหน่วยพื้นที่
\item ค) การเปลี่ยนแปลงความยาว
\item ง) อัตราส่วนระหว่างแรงกับการยืด
\end{enumerate}
\item ความเดียด (Strain) คืออะไร \hfill \textbf{\small (Main)}
\begin{enumerate}[label=)]
\item ก) แรงที่กระทำต่อวัตถุ
\item ข) แรงต่อหน่วยพื้นที่
\item ค) อัตราส่วนระหว่างการเปลี่ยนแปลงความยาวกับความยาวเดิม
\item ง) ความสามารถในการกลับคืนสภาพเดิม
\end{enumerate}
\item มอดูลัสของยัง (Young's Modulus) หาได้จากสูตรใด \hfill \textbf{\small (Main)}
\begin{enumerate}[label=)]
\item ก) Y = F/A
\item ข) Y = (F/A) / (ΔL/L)
\item ค) Y = FL/AΔL
\item ง) Y = ΔL/(FL)
\end{enumerate}
\item พลังงานศักย์ยืดหยุ่น (Elastic Potential Energy) ในสปริงหาได้จากสูตรใด \hfill \textbf{\small (Main)}
\begin{enumerate}[label=)]
\item ก) U = ½kx²
\item ข) U = kx²
\item ค) U = ½kx
\item ง) U = mgx
\end{enumerate}
\item ขีดจำกัดการยืดหยุ่น (Elastic Limit) คืออะไร \hfill \textbf{\small (Main)}
\begin{enumerate}[label=)]
\item ก) จุดที่วัตถุขาด
\item ข) จุดสูงสุดที่วัตถุสามารถยืดแล้วกลับคืนสภาพเดิมได้
\item ค) จุดที่วัตถุเริ่มเสียรูปถาวร
\item ง) จุดที่ความเค้นมากที่สุด
\end{enumerate}
\item เส้นลวดยาว 2 m มีพื้นที่หน้าตัด 1×10⁻⁶ m² ถูกดึงด้วยแรง 100 N ทำให้ยาวขึ้น 0.001 m มอดูลัสของยังมีค่าเท่าไร \hfill \textbf{\small (Main)}
\begin{enumerate}[label=)]
\item ก) 1×10¹¹ Pa
\item ข) 1×10⁸ Pa
\item ค) 1×10¹⁰ Pa
\item ง) 2×10¹¹ Pa
\end{enumerate}
\item เสาเหล็กพื้นที่หน้าตัด 0.01 m² รับแรงกด 50,000 N ความเค้นในเสาเป็นเท่าไร \hfill \textbf{\small (Main)}
\begin{enumerate}[label=)]
\item ก) 5×10⁶ Pa
\item ข) 5×10⁵ Pa
\item ค) 5×10⁴ Pa
\item ง) 5×10⁷ Pa
\end{enumerate}
\item สปริงมีค่า k = 500 N/m ถูกยืดออก 0.1 m พลังงานศักย์ยืดหยุ่นมีค่าเท่าไร \hfill \textbf{\small (Main)}
\begin{enumerate}[label=)]
\item ก) 25 J
\item ข) 2.5 J
\item ค) 50 J
\item ง) 5 J
\end{enumerate}
\item ความเดียด (Strain) มีหน่วยหรือไม่ \hfill \textbf{\small (Main)}
\begin{enumerate}[label=)]
\item ก) มีหน่วยเป็นเมตร
\item ข) มีหน่วยเป็น N/m²
\item ค) ไม่มีหน่วย เป็นปริมาณไร้หน่วย
\item ง) มีหน่วยเป็น Pa
\end{enumerate}
\item วัสดุ A มีมอดูลัสของยังสูงกว่าวัสดุ B หมายความว่าอะไร \hfill \textbf{\small (Main)}
\begin{enumerate}[label=)]
\item ก) วัสดุ A แข็งกว่า แต่เปราะกว่า
\item ข) วัสดุ A ทนต่อแรงดึงได้ดีกว่าต่อการยืดที่เท่ากัน
\item ค) วัสดุ A ยืดได้มากกว่า
\item ง) วัสดุ A หนักกว่า
\end{enumerate}
\item สปริงถูกอัดจากตำแหน่งสมดุล 5 cm ด้วยแรง 20 N พลังงานศักย์ยืดหยุ่นมีค่าเท่าไร \hfill \textbf{\small (Main)}
\begin{enumerate}[label=)]
\item ก) 0.5 J
\item ข) 1 J
\item ค) 2 J
\item ง) 4 J
\end{enumerate}
\item ถ้าแรงดึงเพิ่มขึ้นจนเกินขีดจำกัดการยืดหยุ่น วัสดุจะเป็นอย่างไร \hfill \textbf{\small (Main)}
\begin{enumerate}[label=)]
\item ก) กลับคืนสภาพเดิมทันที
\item ข) เสียรูปถาวรและไม่กลับคืนสภาพเดิม
\item ค) ขาดทันที
\item ง) ไม่เกิดการเปลี่ยนแปลงใดๆ
\end{enumerate}
\item เส้นลวดเหล็กยาว 1 m พื้นที่หน้าตัด 2×10⁻⁶ m² มี Y = 2×10¹¹ Pa ถ้าต้องการยืดให้ยาวขึ้น 0.5 mm ต้องใช้แรงเท่าไร \hfill \textbf{\small (Main)}
\begin{enumerate}[label=)]
\item ก) 200 N
\item ข) 400 N
\item ค) 100 N
\item ง) 800 N
\end{enumerate}
\item เส้นลวดเส้นหนึ่งมีเส้นผ่านศูนย์กลาง 2 mm รับแรงดึง 628 N ความเค้นในเส้นลวดมีค่าเท่าไร \hfill \textbf{\small (Main)}
\begin{enumerate}[label=)]
\item ก) 2×10⁸ Pa
\item ข) 2×10⁷ Pa
\item ค) 2×10⁶ Pa
\item ง) 2×10⁵ Pa
\end{enumerate}
\item แท่งเหล็กยาว 2 m ถูกยืดจนยาว 2.002 m ความเดียดมีค่าเท่าไร \hfill \textbf{\small (Main)}
\begin{enumerate}[label=)]
\item ก) 0.001
\item ข) 0.0001
\item ค) 0.01
\item ง) 0.00001
\end{enumerate}
\item สปริงสองอัน อันที่ 1 มี k = 100 N/m อันที่ 2 มี k = 200 N/m ถูกยืดด้วย x เท่ากัน พลังงานในสปริงอันไหนมากกว่า \hfill \textbf{\small (Main)}
\begin{enumerate}[label=)]
\item ก) สปริงที่ 1
\item ข) สปริงที่ 2
\item ค) เท่ากัน
\item ง) ขึ้นกับทิศทาง
\end{enumerate}
\item มอดูลัสของยังของเหล็ก ≈ 2×10¹¹ Pa และของอะลูมิเนียม ≈ 0.7×10¹¹ Pa ข้อสรุปใดถูกต้อง \hfill \textbf{\small (Main)}
\begin{enumerate}[label=)]
\item ก) เหล็กยืดได้ง่ายกว่าอะลูมิเนียม
\item ข) เหล็กแข็งแกร่งกว่าอะลูมิเนียม
\item ค) อะลูมิเนียมทนต่อแรงดึงได้ดีกว่า
\item ง) เหล็กหนักกว่า
\end{enumerate}
\item จุดปลายยืดหยุ่น (Proportional Limit) แตกต่างจากขีดจำกัดการยืดหยุ่น (Elastic Limit) อย่างไร \hfill \textbf{\small (Main)}
\begin{enumerate}[label=)]
\item ก) จุดปลายยืดหยุ่นอยู่หลังขีดจำกัดการยืดหยุ่น
\item ข) จุดปลายยืดหยุ่นอยู่ก่อนขีดจำกัดการยืดหยุ่น
\item ค) ทั้งสองคือจุดเดียวกัน
\item ง) ไม่มีความแตกต่าง
\end{enumerate}
\item สปริงมี k = 800 N/m ถูกยืดจากตำแหน่งสมดุลไป 10 cm แล้วยืดต่ออีก 10 cm รวม 20 cm พลังงานเพิ่มขึ้นจากเท่าไร \hfill \textbf{\small (Main)}
\begin{enumerate}[label=)]
\item ก) 0.4 J → 1.6 J เพิ่ม 1.2 J
\item ข) 0.4 J → 1.6 J เพิ่ม 4 J
\item ค) 0.8 J → 3.2 J เพิ่ม 2.4 J
\item ง) 0.4 J → 0.8 J เพิ่ม 0.4 J
\end{enumerate}
\item เมื่อวัสดุถูกกด (Compression) ความเค้นที่เกิดขึ้นเรียกว่าอะไร \hfill \textbf{\small (Main)}
\begin{enumerate}[label=)]
\item ก) Compressive stress
\item ข) Tensile stress
\item ค) Shear stress
\item ง) Bulk stress
\end{enumerate}
\item เสาคอนกรีตสี่เหลี่ยมหน้ากว้าง 0.3 m × 0.3 m สูง 3 m รับน้ำหนัก 5×10⁵ N ความเค้นกดในเสาเป็นเท่าไร \hfill \textbf{\small (Main)}
\begin{enumerate}[label=)]
\item ก) ≈ 5.6×10⁶ Pa
\item ข) ≈ 1.7×10⁶ Pa
\item ค) ≈ 5.6×10⁵ Pa
\item ง) ≈ 1.7×10⁷ Pa
\end{enumerate}
\item สายเคเบิลยาว 10 m ถูกยืดจนยาว 10.01 m ความเดียดมีค่าเท่าไร \hfill \textbf{\small (Main)}
\begin{enumerate}[label=)]
\item ก) 0.001
\item ข) 0.0001
\item ค) 0.01
\item ง) 0.1
\end{enumerate}
\item สปริงมี k = 200 N/m ถูกยืด 0.05 m จากตำแหน่งสมดุล แล้วปล่อย สปริงจะมีพลังงานจลน์สูงสุดเท่าไรที่ตำแหน่งสมดุล \hfill \textbf{\small (Main)}
\begin{enumerate}[label=)]
\item ก) 0.5 J
\item ข) 0.25 J
\item ค) 0.125 J
\item ง) 0.1 J
\end{enumerate}
\item ลวดทองแดงยาว 4 m พื้นที่หน้าตัด 2×10⁻⁶ m² มี Y = 1.1×10¹¹ Pa ต้องใช้แรงเท่าไรเพื่อยืดให้ยาวขึ้น 1 mm \hfill \textbf{\small (Main)}
\begin{enumerate}[label=)]
\item ก) 55 N
\item ข) 110 N
\item ค) 220 N
\item ง) 55,000 N
\end{enumerate}
\item กฎของฮุค (Hooke's Law) ระบุว่าอะไร \hfill \textbf{\small (Main)}
\begin{enumerate}[label=)]
\item ก) F = -kx
\item ข) σ = Yε
\item ค) ทั้ง ก และ ข ถูก
\item ง) F = kx²
\end{enumerate}
\item ความเค้น (Stress) มีหน่วยเป็นอะไร \hfill \textbf{\small (Extra)}
\begin{enumerate}[label=)]
\item ก) N
\item ข) m
\item ค) Pa หรือ N/m²
\item ง) J
\end{enumerate}
\item มอดูลัสของยังของเหล็ก ≈ 2×10¹¹ Pa หมายถึงอะไร \hfill \textbf{\small (Extra)}
\begin{enumerate}[label=)]
\item ก) เหล็กยืดได้ 2×10¹¹ เมตร
\item ข) เหล็กทนต่อแรงได้ดีมาก โดยต้องใช้แรงมากมายเพื่อให้เกิดการยืดเล็กน้อย
\item ค) เหล็กมีความหนาแน่นสูง
\item ง) เหล็กขยายตัวได้ 2×10¹¹ เท่า
\end{enumerate}
\item ความเดียด (Strain) คำนวณจากอะไร \hfill \textbf{\small (Extra)}
\begin{enumerate}[label=)]
\item ก) ΔL × L₀
\item ข) ΔL / L₀
\item ค) L₀ / ΔL
\item ง) F / A
\end{enumerate}
\item สปริงถูกยืด 0.02 m จากตำแหน่งสมดุล มีพลังงานศักย์ 0.5 J ค่า k ของสปริงเป็นเท่าไร \hfill \textbf{\small (Extra)}
\begin{enumerate}[label=)]
\item ก) 250 N/m
\item ข) 500 N/m
\item ค) 1,000 N/m
\item ง) 2,500 N/m
\end{enumerate}
\item วัสดุที่มีขีดจำกัดการยืดหยุ่นสูง หมายถึงอะไร \hfill \textbf{\small (Extra)}
\begin{enumerate}[label=)]
\item ก) วัสดุนั้นแข็งมาก
\item ข) วัสดุนั้นสามารถรับแรงได้มากก่อนจะเกิดการเสียรูปถาวร
\item ค) วัสดุนั้นยืดได้มาก
\item ง) วัสดุนั้นเปราะ
\end{enumerate}
\item แท่งอะลูมิเนียมยาว 1 m พื้นที่หน้าตัด 1×10⁻⁴ m² รับแรงดึง 2,000 N ยืด 0.5 mm มอดูลัสของยังเป็นเท่าไร \hfill \textbf{\small (Extra)}
\begin{enumerate}[label=)]
\item ก) 4×10¹⁰ Pa
\item ข) 2×10¹⁰ Pa
\item ค) 4×10⁹ Pa
\item ง) 2×10¹¹ Pa
\end{enumerate}
\item สายเชือกเส้นหนึ่งมีเส้นผ่านศูนย์กลาง 4 mm รับแรงดึง 200 N ความเค้นในเชือกเป็นเท่าไร \hfill \textbf{\small (Extra)}
\begin{enumerate}[label=)]
\item ก) 1.6×10⁷ Pa
\item ข) 1.6×10⁶ Pa
\item ค) 1.6×10⁵ Pa
\item ง) 6.4×10⁷ Pa
\end{enumerate}
\item สปริงถูกยืดจาก 5 cm เป็น 10 cm จากตำแหน่งสมดุล ถ้า k = 400 N/m พลังงานศักย์สุดท้ายเป็นเท่าไร \hfill \textbf{\small (Extra)}
\begin{enumerate}[label=)]
\item ก) 0.5 J
\item ข) 2 J
\item ค) 0.2 J
\item ง) 4 J
\end{enumerate}
\item ภายในขีดจำกัดการยืดหยุ่น ความสัมพันธ์ระหว่างความเค้นและความเดียดเป็นอย่างไร \hfill \textbf{\small (Extra)}
\begin{enumerate}[label=)]
\item ก) σ ∝ ε²
\item ข) σ = Y/ε
\item ค) σ ∝ ε
\item ง) σ = คงที่
\end{enumerate}
\item เสาอลูมิเนียมพื้นที่หน้าตัด 0.001 m² สูง 5 m มี Y = 0.7×10¹¹ Pa ถ้าถูกกดด้วยแรง 3.5×10⁶ N จะหดสั้นลงเท่าไร \hfill \textbf{\small (Extra)}
\begin{enumerate}[label=)]
\item ก) 0.025 m
\item ข) 0.25 m
\item ค) 0.0025 m
\item ง) 0.05 m
\end{enumerate}
\item ความเดียด 0.002 หมายความว่าอะไร \hfill \textbf{\small (Extra)}
\begin{enumerate}[label=)]
\item ก) ความยาวเพิ่มขึ้น 0.2\%
\item ข) ความยาวเพิ่มขึ้น 0.02\%
\item ค) ความยาวเพิ่มขึ้น 2\%
\item ง) ความยาวเพิ่มขึ้น 20\%
\end{enumerate}
\item สปริง k = 1,000 N/m ถูกอัดจากความยาวธรรมชาติ 0.3 m เป็น 0.2 m พลังงานศักย์เป็นเท่าไร \hfill \textbf{\small (Extra)}
\begin{enumerate}[label=)]
\item ก) 5 J
\item ข) 50 J
\item ค) 0.5 J
\item ง) 0.05 J
\end{enumerate}
\item ทองแดงมี Y = 1.1×10¹¹ Pa และเหล็กมี Y = 2×10¹¹ Pa ถ้าใช้แรงเท่ากันดึงเส้นลวดทั้งสองที่มีพื้นที่หน้าตัดเท่ากัน ข้อใดถูกต้อง \hfill \textbf{\small (Extra)}
\begin{enumerate}[label=)]
\item ก) เหล็กยืดน้อยกว่า
\item ข) ทองแดงยืดน้อยกว่า
\item ค) ทั้งสองยืดเท่ากัน
\item ง) ไม่สามารถระบุได้
\end{enumerate}
\item เมื่อแรงดึงเพิ่มขึ้น 2 เท่า แต่พื้นที่หน้าตัดคงที่ ความเค้นจะเป็นอย่างไร \hfill \textbf{\small (Extra)}
\begin{enumerate}[label=)]
\item ก) เพิ่มขึ้น 2 เท่า
\item ข) ลดลง 2 เท่า
\item ค) คงที่
\item ง) เพิ่มขึ้น 4 เท่า
\end{enumerate}
\item จุดที่วัสดุเริ่มขาดเรียกว่าอะไร \hfill \textbf{\small (Extra)}
\begin{enumerate}[label=)]
\item ก) Elastic limit
\item ข) Proportional limit
\item ค) Breaking point
\item ง) Yield point
\end{enumerate}
\item สปริงสองอันที่มี k = 200 N/m และ k = 400 N/m ต่ออนุกรมกัน ค่า k รวมเป็นเท่าไร \hfill \textbf{\small (Extra)}
\begin{enumerate}[label=)]
\item ก) 600 N/m
\item ข) 133 N/m
\item ค) 800 N/m
\item ง) 200 N/m
\end{enumerate}
\item สายเคเบิลเหล็กยาว 10 m รัศมี 0.005 m มี Y = 2×10¹¹ Pa ถ้าต้องการยืดไม่เกิน 0.002 m จะรับแรงได้สูงสุดเท่าไร \hfill \textbf{\small (Extra)}
\begin{enumerate}[label=)]
\item ก) ≈ 1.57×10⁴ N
\item ข) ≈ 3.14×10⁴ N
\item ค) ≈ 7.85×10⁴ N
\item ง) ≈ 15.7×10⁴ N
\end{enumerate}
\item แท่งเหล็กยาว 5 m ถูกยืดจนมีความเดียด 0.003 ความยาวสุดท้ายเป็นเท่าไร \hfill \textbf{\small (Extra)}
\begin{enumerate}[label=)]
\item ก) 5.015 m
\item ข) 5.003 m
\item ค) 5.03 m
\item ง) 5.3 m
\end{enumerate}
\item สปริง k = 300 N/m ถูกยืด 0.1 m แล้วปล่อย ไม่มีแรงต้าน สปริงจะมีพลังงานจลน์สูงสุดเท่าไร \hfill \textbf{\small (Extra)}
\begin{enumerate}[label=)]
\item ก) 1.5 J
\item ข) 3 J
\item ค) 0.75 J
\item ง) 6 J
\end{enumerate}
\item เสาประตูมีพื้นที่หน้าตัด 0.005 m² รับแรงกด 10,000 N ความเค้นกดมีค่าเท่าใด \hfill \textbf{\small (Extra)}
\begin{enumerate}[label=)]
\item ก) 2×10⁶ Pa
\item ข) 2×10⁵ Pa
\item ค) 2×10⁷ Pa
\item ง) 2×10⁴ Pa
\end{enumerate}
\item ค่ามอดูลัสของยังของวัสดุขึ้นกับอะไร \hfill \textbf{\small (Extra)}
\begin{enumerate}[label=)]
\item ก) ขึ้นกับขนาดของวัสดุ
\item ข) ขึ้นกับชนิดของวัสดุ ไม่ขึ้นกับขนาด
\item ค) ขึ้นกับแรงที่กระทำ
\item ง) ขึ้นกับการยืดที่เกิดขึ้น
\end{enumerate}
\item ข้อใดไม่ใช่คุณสมบัติของขีดจำกัดการยืดหยุ่น \hfill \textbf{\small (Extra)}
\begin{enumerate}[label=)]
\item ก) เป็นจุดสูงสุดที่วัสดุยังกลับคืนสภาพเดิมได้
\item ข) เป็นค่าคงที่ของวัสดุ
\item ค) เป็นจุดที่วัสดุขาด
\item ง) เกินจุดนี้วัสดุจะเสียรูปถาวร
\end{enumerate}
\item สปริงสองอัน k = 100 N/m ต่อขนานกัน ค่า k รวมเป็นเท่าไร \hfill \textbf{\small (Extra)}
\begin{enumerate}[label=)]
\item ก) 100 N/m
\item ข) 200 N/m
\item ค) 50 N/m
\item ง) 10 N/m
\end{enumerate}
\item ความเค้นดึง (Tensile stress) เกิดขึ้นเมื่อใด \hfill \textbf{\small (Extra)}
\begin{enumerate}[label=)]
\item ก) เมื่อแรงกดทำให้วัสดุสั้นลง
\item ข) เมื่อแรงดึงทำให้วัสดุยาวขึ้น
\item ค) เมื่อแรงเฉือนทำให้วัสดุบิดเบี้ยว
\item ง) เมื่อแรงดันทำให้วัสดุพอง
\end{enumerate}
\item ลวดทองเทียมยาว L พื้นที่ A มีมอดูลัสของยัง Y ถูกดึงด้วยแรง F จนยาวขึ้น ΔL ข้อใดถูกต้อง \hfill \textbf{\small (Extra)}
\begin{enumerate}[label=)]
\item ก) ΔL = FLY/A
\item ข) ΔL = FL/(AY)
\item ค) ΔL = FAY/L
\item ง) ΔL = YAF/L
\end{enumerate}
\item สปริง k = 500 N/m ถูกยืดจากตำแหน่งสมดุลไป 10 cm แล้วปล่อย สปริงจะเดินทางผ่านตำแหน่งสมดุลไปได้อีกกี่ cm ก่อนจะหยุด (ไม่คิดแรงต้าน) \hfill \textbf{\small (Extra)}
\begin{enumerate}[label=)]
\item ก) 10 cm
\item ข) 5 cm
\item ค) 20 cm
\item ง) 0 cm
\end{enumerate}
\item วัสดุที่มีความเดียดสูง หมายถึงอะไร \hfill \textbf{\small (Extra)}
\begin{enumerate}[label=)]
\item ก) วัสดุนั้นแข็งแกร่งมาก
\item ข) วัสดุนั้นยืดได้มากภายใต้แรงเท่ากัน
\item ค) วัสดุนั้นเปราะ
\item ง) วัสดุนั้นหนัก
\end{enumerate}
\item เหล็กกล้ามีมอดูลัสของยัง 200 GPa มีค่าเท่ากับกี่ Pa \hfill \textbf{\small (Extra)}
\begin{enumerate}[label=)]
\item ก) 2×10¹¹ Pa
\item ข) 2×10⁹ Pa
\item ค) 2×10¹⁰ Pa
\item ง) 2×10⁸ Pa
\end{enumerate}
\item ทำไมสปริงจึงกลับคืนสภาพเดิมได้หลังถูกดึงหรืออัด \hfill \textbf{\small (Extra)}
\begin{enumerate}[label=)]
\item ก) เพราะสปริงมีมวลมาก
\item ข) เพราะสปริงทำจากโลหะที่มีขีดจำกัดการยืดหยุ่นสูง
\item ค) เพราะสปริงมีพลังงานจลน์
\item ง) เพราะสปริงไม่เกิดความเดียด
\end{enumerate}
\item สปริง k = 200 N/m ต่ออนุกรมกับสปริง k = 300 N/m ถูกยืด 0.15 m พลังงานศักย์รวมเป็นเท่าไร \hfill \textbf{\small (Extra)}
\begin{enumerate}[label=)]
\item ก) ≈ 1.35 J
\item ข) ≈ 2.25 J
\item ค) ≈ 3.75 J
\item ง) ≈ 4.5 J
\end{enumerate}
\end{enumerate}
\end{document}
